\section{The retained cover theorem and its exact residue transport}\label{NC.residues}
Throughout this supplement the complete common paper~\cite{OldCommon}
is an acquired dependency. In particular Theorem C.Graph proves the
geometry, thinness, natural faithful chamber action, and exact
residual-connectivity criterion of the component construction. Its proof
includes the connected-truncation lemma, the component-intersection
argument and extension of every flag. The following fuller residue
identification is the accepted phase27 argument~\cite{Parity}; it is
needed to recover residue diagrams rather than connectivity alone.

Let $\Gamma$ be a residually connected incidence geometry of finite rank
at least three, satisfying (B1) and (B2),
with connected simple graph $Q$ on types $0,1$. For a flag $F$ of higher
objects write $Q_F$ for its common vertex--edge shadow. Connectedness of
these graphs follows by applying the connected-truncation lemma inside
the original residues, and chamber extension supplies an edge. Write
$D(S)$ for the graph with vertices $(v,\epsilon)$ and edges
$(v,0)(w,1)$ whenever $vw$ is an edge of $S$. Choose a component $C$
of $D(Q)$. A higher object of the transformed system is $(x,C_x)$,
where $C_x$ is a component of $D(Q_x)$ contained in $C$.

This is literally Construction 3.6 in~\cite{PT} when $Q$ is
nonbipartite: the type-$0$ part of $C_x$ is an even-walk class, and
its type-$1$ part is the adjacent class, or the same full shadow in the
nonbipartite case. If two such components meet at type $1$, a neighbor
in the connected common graph $Q_{x,y}$ gives a common type-$0$ point.
Thus the stated component incidence is exactly the source incidence,
not merely a bijection on its objects. When $Q$ is bipartite, each
$D(Q_x)$ has exactly one component inside $C$; this is its
Construction 3.1, attributed there to Lef\`evre-Percsy--Percsy--Leemans.

For a higher flag $\Phi=\{(x,C_x):x\in F\}$, put
$W_\Phi=\bigcap_{x\in F}C_x$, with $W_\varnothing=C$.
The component lemma of C.Graph states that $W_\Phi$ is one component
of $D(Q_F)$ if any $Q_x$ is bipartite, and is all of $D(Q_F)$
otherwise. Its connectedness is therefore governed by
\[
 \tag{PC}
 [Q_x\text{ nonbipartite for every }x\in F]
       \ \Longrightarrow\ Q_F\text{ nonbipartite}.
\]
The condition ranges over actual higher flags, not arbitrary sets of
objects. It is automatic for globally bipartite $Q$.

\begin{theorem}[Full residue transport]\label{NC.transport}
Assume (PC). For a flag of higher lifts $\Phi$, its residue is
isomorphic to the component construction on the original residue
$\Gamma_F$, with selected component $W_\Phi$ of $D(Q_F)$.
If a defining flag contains a new vertex, its residue is isomorphic
to the corresponding original residue. The local component construction
again satisfies its parity-compatibility requirement.
\end{theorem}
\begin{proof}
The vertex case is the explicit isomorphism of C.Graph: the new
neighbor of $(v,0)$ is replaced by its unique original edge at $v$,
and a higher lift is replaced by its original object. Its component
is uniquely determined by $(v,0)$. Simplicity and (B2) preserve
incidence. Restrict this map for a larger flag containing that vertex.

For a flag of higher lifts, (PC) makes $W_\Phi$ connected. Suppose
first that $Q_F$ is bipartite. For any further object $y$ incident
with $F$, its connected shadow $Q_{F,y}$ has exactly one component
of its double cover in $W_\Phi$. That component lies in a unique
component of $D(Q_y)$. This component meets each prescribed $C_x$
over $Q_{F,y}$, and therefore gives the unique compatible global
lift. Conversely a compatible global lift restricts to that component:
apply the component-intersection lemma to $F\cup\{y\}$. This
establishes the object bijection, including when $F$ is empty.

If $Q_F$ is nonbipartite, every $Q_x$, $x\in F$, is nonbipartite.
For each further $y$, (PC) on $F\cup\{y\}$ shows that $Q_{F,y}$
is nonbipartite exactly when $Q_y$ is nonbipartite. Indeed the forward
implication follows from containment, and the reverse is precisely
(PC). Thus restriction from components of $D(Q_y)$ to components
of $D(Q_{F,y})$ is a bijection, with either one or two components.
Each restriction is compatible with every member of $\Phi$.

In either case, apply the same component-intersection argument to
$F\cup\{y,z\}$. Compatible components for an incident pair have
nonempty common restriction exactly when their corresponding global
lifts are incident. Here a prescribed bipartite constituent selects
the same component on every connected subshadow; if all constituents
are nonbipartite, (PC) makes the total shadow nonbipartite. This proves
preservation of higher incidence. Vertex incidence is membership, so
is preserved directly. These maps identify the complete residues.

Finally, in the nonbipartite local case, if every singleton shadow
$Q_{F,y}$ of a local flag is nonbipartite, all $Q_y$ and all $Q_x$
for $x\in F$ are nonbipartite. Applying global (PC) to their union
makes their total common shadow nonbipartite. The bipartite local
case uses the bipartition construction and needs no further condition.
\end{proof}

In a thin rank-two residue, a finite polygon shadow of length $m$
gives a cover component of length $m$ when $m$ is even and $2m$
when $m$ is odd. Its new gonality is therefore $m/2$ or $m$.
For a bi-infinite polygon both components are again bi-infinite paths.
The vertex-residue isomorphisms duplicate the adjacent diagram labels.
Thus the actual rank-two residues, including infinite ones, are
accounted for. Without (PC), the component-intersection description
still holds and the transformed system is still a geometry, but the
stronger single-component residue assertion is not available.

\subsection{Relator parity and the regular-leaf hypothesis}
The elementary group calculation $H_J\trianglelefteq U_J=H_J\gen a$
in C.Endpoint also proves the word form of source Proposition 4.3.
For a presentation on $a,b,r_i$, assign $1\in\F_2$ to $a$ and
$0$ to every other marker. This gives a homomorphism from the free
group. It factors through the presented group exactly when every
relator has even total $a$-parity, since then the relators and their
normal closure lie in its kernel. If it factors, its kernel contains
$\gen{aba,b,r_i}$ and excludes $a$, so the index bound identifies
that kernel with the halving subgroup. Conversely, if that subgroup
is proper, its quotient character has these generator values, and
every relator must have even $a$-parity. The parity lemma below
identifies this condition with bipartiteness of the vertex graph.
This proves all four equivalences with a fixed presentation; no
finite order computation is used.

For completeness, the commuting-marker hypothesis at a regular leaf
is supplied by the geometry. Fix all types other than $0,i$, with
$i\ge2$. Its fixed type-$1$ edge has two endpoints. (B2) says that
any compatible type-$i$ object is incident with both. Thinness makes
this rank-two residue a digon. Its two panel involutions therefore
commute, giving $[a,r_i]=1$. This is the exact bridge from (B1)/(B2)
and the regular input to the algebra used below.
